% Part: first-order-logic
% Chapter: syntax-and-semantics
% Section: subformulas

\documentclass[../../../include/open-logic-section]{subfiles}

\begin{document}

\olfileid[pt]{fol}{syn}{sbf}

\olsection{\printtoken{P}{subformula}}

\begin{explain}
Frequentemente é útil falar sobre as !!{formula}s que ``compõem'' uma
dada !!{formula}.  Chamamos estas de suas \emph{!!{subformula}s}.  Qualquer
!!{formula} conta como !!a{subformula} de si mesma; uma subfórmula de $!A$
diferente do próprio $!A$ é uma \emph{!!{subformula} própria}.
\end{explain}

\begin{defn}[!!^{subformula} imediata]
Se $!A$ é !!a{formula}, as \emph{!!{subformula}s imediatas}
de $!A$ são definidas indutivamente como segue:
\begin{enumerate}
\item !!^{formula}s atômicas não têm !!{subformula}s imediatas.

\tagitem{prvNot}{\indcase{!A}{\lnot !B}{A única
    !!{subformula} imediata de $\indfrm$ é~$!B$.}}{}

\item \indcase{!A}{(!B \ast !C)}{As !!{subformula}s imediatas de
  $\indfrm$ são $!B$ e $!C$ ($\ast$ é qualquer um dos conectivos
  de dois lugares).}

\tagitem{prvAll}{\indcase{!A}{\lforall[x][!B]}{A única
    !!{subformula} imediata de $\indfrm$ é~$!B$.}}{}

\tagitem{prvEx}{\indcase{!A}{\lexists[x][!B]}{A única
    !!{subformula} imediata de $\indfrm$ é~$!B$.}}{}
\end{enumerate}
\end{defn}

\begin{defn}[!!^{subformula} própria]
Se $!A$ é !!a{formula}, as \emph{!!{subformula}s próprias}
de $!A$ são definidas recursivamente como segue:
\begin{enumerate}
\item !!^{formula}s atômicas não têm !!{subformula}s próprias.

\tagitem{prvNot}{\indcase{!A}{\lnot !B}{As !!{subformula}s próprias de
    $\indfrm$ são~$!B$ juntamente com todas as !!{subformula}s próprias
    de~$!B$.}}{}

\item \indcase{!A}{(!B \ast !C)}{As !!{subformula}s próprias de
  $\indfrm$ são $!B$, $!C$, juntamente com todas as !!{subformula}s próprias
  de $!B$ e as de~$!C$.}

\tagitem{prvAll}{\indcase{!A}{\lforall[x][!B]}{As !!{subformula}s
    próprias de $\indfrm$ são~$!B$ juntamente com todas as
    !!{subformula}s próprias de~$!B$.}}{}

\tagitem{prvEx}{\indcase{!A}{\lexists[x][!B]}{As !!{subformula}s
    próprias de $\indfrm$ são~$!B$ juntamente com todas as
    !!{subformula}s próprias de~$!B$.}}{}
\end{enumerate}
\end{defn}

\begin{defn}[!!^{subformula}]
As !!{subformula}s de $!A$ são o próprio $!A$ juntamente com todas as suas
!!{subformula}s próprias.
\end{defn}

\begin{explain}
Note a diferença sutil em como definimos !!{subformula}s
imediatas e !!{subformula}s próprias.  No primeiro caso,
definimos diretamente as !!{subformula}s imediatas de uma fórmula~$!A$
para cada forma possível de~$!A$.  É uma definição explícita por
casos, e os casos espelham a definição indutiva do conjunto de
!!{formula}s.  No segundo caso, também espelhamos a maneira como o
conjunto de todas as !!{formula}s é definido, mas em cada caso também
incluímos as !!{subformula}s próprias das !!{formula}s menores $!B$,
$!C$ além dessas próprias !!{formula}s.  Isso torna a
definição \emph{recursiva}.  Em geral, uma definição de uma função
sobre um conjunto definido indutivamente (no nosso caso, !!{formula}s) é recursiva
se os casos na definição da função fazem uso da
própria função. Para ser bem definida, devemos garantir, no entanto,
que só usamos os valores da função para argumentos que
vêm ``antes'' daquele que estamos definindo---no nosso caso, ao definir
``!!{subformula} própria'' para $(!B \ast !C)$ usamos apenas as
!!{subformula}s próprias das !!{formula}s ``anteriores'' $!B$ e $!C$.
\end{explain}

\begin{prop}
\ollabel{prop:subfrm-trans}
Suponha que $!B$ é uma subfórmula de $!A$ e $!C$ é uma subfórmula de $!B$.
Então $!C$ é uma subfórmula de $!A$. Em outras palavras, a relação de subfórmula
é transitiva.
\end{prop}

\begin{prob}
Prove \olref[fol][syn][sbf]{prop:subfrm-trans}.
\end{prob}

\begin{prop}
\ollabel{prop:count-subfrms}
Suponha que $!A$ é uma fórmula com $n$ conectivos e quantificadores.
Então $!A$ tem no máximo $2n+1$ subfórmulas.
\end{prop}

\begin{prob}
Prove \olref[fol][syn][sbf]{prop:count-subfrms}.
\end{prob}

\end{document}
